\section{Style measurement and final review}\label{reward-adjudication}

This section defines the style utility \(V\) and semantic gate \(G\)
for B transformation. Every B task requires preservation of the source's
full meaning, fulfillment of the task's target tonal intent, and
readability. These constraints
create no separate writing-improvement service or reward pool. Positive
utility requires independent semantic certification; style proximity or
detector evasion cannot compensate for a preservation failure.
Numerical measurement uses a public frozen artifact.
Semantic judgments are explicit oracle inputs certified under the
validator-weight assumption in Section~\ref{service-evidence}.

Validators obtain authoritative benchmark origin probabilities by
executing the published A panel under the artifact-freeze and execution
rules in Sections~\ref{a-artifact-execution}
and~\ref{common-transformation-evidence}. They
compare B's submitted vectors with canonical outputs bound to the same
artifact, runtime and input. B models remain private; rewards assess
their committed output text. B miners can compute public A scores during
search, while hidden labels and unreleased semantic-review evidence
remain restricted.

\paragraph{Frozen style measurement.}
Style measurement applies to tasks requesting target-author fit.
Evasion-only tasks record it as \texttt{not\_requested} and need no
target-reference or calibration input.
Before task issue, publish a permitted word/grammar author-distance
artifact and pin its hash, parser and feature schemas, model parameters,
reference-pooling rule, executable/runtime, numerical operations,
rounding, resource bounds, and test vectors. Its distance function
\(d_\theta(t,R)\) returns a nonnegative integer for text \(t\) and
target-author references \(R\). The instance manifest must supply these
concrete bindings; this specification does not select or claim to have
trained the artifact. Private B generator weights do not define this
public style evaluator.

Freeze a nonempty calibration multiset \(C=(c_1,\ldots,c_M)\) of
integer distances from development queries to their attributed authors'
independent reference works. Keep development source families separate
from fitting and evaluation, record provenance and hashes, and use the
same artifact and reference-pooling rule. The task pins its calibration
ID before candidate generation. Retain duplicate distances. With
\(Q=1{,}000{,}000\), define the integer style fit
\begin{equation}\label{eq:style-fit}
 q(t,R)=\left\lfloor
 Q\,\frac{2\#\{m:c_m>d_\theta(t,R)\}
             +\#\{m:c_m=d_\theta(t,R)\}}{2M}
 \right\rfloor.
\end{equation}
This empirical rank gives exact ties half credit. Smaller distance
cannot reduce \(q\). It is a calibrated distance rank, not an authorship
probability or a certificate of preserved meaning.

For source \(x\) and candidate \(z\), evaluate both against the same
\(R\), artifact, and calibration. Let \(q_x=q(x,R)\), \(q_z=q(z,R)\),
and define
\begin{equation}\label{eq:style-utility}
 V(x,z)=\begin{cases}
 0,&q_x=Q,\\
 \displaystyle\frac{\max(0,q_z-q_x)}{Q-q_x},&q_x<Q.
 \end{cases}
\end{equation}
Equal fit earns zero improvement; a saturated source has no remaining
headroom. Preserve the exact rational numerator and denominator through
aggregation. A task requesting no style regression requires
\(q_z\geq q_x\), even though \(V\) clips negative improvement.
Task mixtures and final credit rounding remain part of the frozen
allocation policy.

For style-scored tasks, source and reference scoring must succeed before
issue. A declared, deterministic \texttt{unscorable\_candidate} result gives zero style
utility. Execution disagreement or unavailable shared evidence cannot
silently supply a score: it remains unresolved under the common
comparison-closure rule. A numerical failure alone does not establish
attributable service nonresponse. Validate the frozen artifact against
copying, topic substitution, and meaning-breaking controls before
funds depend on it.

Miners can also optimize directly against the public style artifact.
Before live rewards, test adaptive searches that maximize \(V\) under
Section~\ref{development-and-launch-criteria}. Among revisions that pass
\(P,R,T\), higher \(V\) must still track independent reader judgments
of author fit; the semantic gates alone do not validate that ranking.

\paragraph{Three semantic fields.}
The source's full meaning must remain intact. Before generation, resolve
the target tonal intent from the brief and, where supplied, the requested
authorial style and authorized reference samples; freeze it in the task rubric. A dramatic
change of authorial style can authorize a different tone without a
separate tone instruction. Use source tone as the default for aspects the
request leaves unchanged. Review the whole candidate for meaning,
including information not enumerated in the checklist. Tone review
applies the frozen rubric and its source-tone defaults throughout the
candidate; reviewers cannot add new stylistic preferences after
generation. Every certifying validator records one
verdict, \texttt{PASS}, \texttt{FAIL}, or \texttt{ABSTAIN}, for each field:
\begin{itemize}
\tightlist
\item
  \(P\), preservation: retain the source's full meaning, including all
  claims, entities, quantities, negation, uncertainty, attribution,
  citations, and argument relations. Changes to expression must preserve
  that information. Unsupported material additions, omissions, altered
  meaning, or prohibited reference copying fail. Greater stylistic force
  cannot strengthen factual certainty or obligations.
\item
  \(R\), readability: introduce no material readability regression
  against the source and satisfy the brief's stated readability
  requirements, including required cleanup of identified unwanted
  filler, repetition or formulaic phrasing. Shorter text or simpler
  vocabulary alone supplies no pass.
\item
  \(T\), tone: fulfill the frozen target tonal intent and every tone
  and audience constraint. Preserve the intended voice features,
  including informality or idiosyncrasies, specified by that target.
  This field always requires a semantic vote, including when no separate
  tone instruction appears in the request.
\end{itemize}
Insufficient evidence or unresolved interpretation, including uncertainty
about the target tonal intent, requires abstention.
Failure ballots identify rubric items and bound source/candidate spans
where applicable. Preference ties and extra editorial praise earn no
additional utility. No detector probability substitutes for these
semantic fields.

\paragraph{Authority and timing.}
Three hash-assigned preparers assemble the dossier; their vote count
cannot finalize a judgment. Every final signer receives authorized
benchmark evidence, checks the relevant source, references, brief,
candidate and rubric, and attests the semantic field itself. The review
client presents semantic evidence first; honest signers commit their
field ballots before inspecting detector scores or numerical style
results. Full-set authorization is not a cryptographic blinding guarantee.
Checking
three preparers' signatures or tally is insufficient. Reuse the verified
frozen validator snapshot and integer effective weights from
Section~\ref{service-evidence}, including its conditional assumption of
strictly less than \(W/3\) dishonest eligible weight. Positive \(W\),
qualified recipient keys and complete capacity reservations are required
before issue.

For the pilot epoch starting at \(H\), preparers commit by \(H+180\)
and open privately by \(H+192\). Evidence corrections close at
\(H+196\); they may correct references to already committed evidence,
but cannot change the candidate, brief, evaluator, or calibration.
Each validator then commits its one final ballot per field by
\(H+198\), opens it privately by \(H+202\), and the global field
certificate must be included by \(H+204\). Deadline equality is timely;
finalized inclusion, incident scope and payout timing follow
Section~\ref{service-evidence}. Reserve all witness publication and
verification work before issue. Under the pilot response schedule,
candidate certification occurs by \(H+72\) through \(H+168\), leaving
30 to 126 blocks before the ballot commitment cutoff: approximately
6 to 25 minutes at 12 seconds per block. Certifiers must review the
authorized source and candidate as evidence becomes available; waiting
for the preparers' private opening would leave only six blocks.
These cutoffs do not establish measured review throughput.

The first canonical timely commitment fixes that validator's ballot for
the field. Identical retries are idempotent; a conflicting signed ballot
is retained as evidence and cannot replace it. Late, absent, or invalid
openings contribute no agreeing weight. Keep missing, abstaining,
recused, and nonresponsive validators in the frozen denominator \(W\).
A \texttt{PASS} or \texttt{FAIL} certificate requires distinct eligible
signers of that verdict with total weight \(w\) satisfying
\(3w>2W\). Honest signers never authorize conflicting final verdicts;
conflicting strict certificates halt affected settlement.

\paragraph{Deterministic closure.}
At the close, set \(G=1\) only if all three fields have a valid
\texttt{PASS} certificate. Set \(G=0\) if any field has a valid
\texttt{FAIL} certificate. Otherwise the judgment is unresolved.
Conflicting certificates take precedence and
halt settlement. Missing candidate submissions remain attributable
miner nonresponse under the service-evidence rules.

An unresolved candidate judgment voids that matched source/brief/mode
comparison for the entire B batch. It cannot selectively remove a
difficult candidate or erase independently attributable service failures.
Retain all ballots and disagreements; do not draw replacement judges
to obtain a preferred verdict. This rule makes outcomes reproducible
from certified inputs. It does not require honest readers to interpret
every passage identically or turn quorum agreement into semantic proof.

Whole-batch closure prevents validators from using unresolved judgments
to exclude only one competitor's candidate, at the cost of giving B
submitters a way to disrupt the comparison. A chosen candidate can split
honest judgments
or induce abstention so that neither \texttt{PASS} nor \texttt{FAIL}
exceeds \(2W/3\) on a required field. If no other field receives a
\texttt{FAIL} certificate, that candidate voids the comparison. Dissent
alone need not cause a void: enough agreeing failure weight certifies
\(G=0\). A dishonest minority below \(W/3\) also cannot prevent a
certificate when all remaining weight agrees and participates.
An unresolved semantic judgment adds no service failure or recovery
deficit; the submitter still pays its generation and any report costs. A UID
expecting zero skill could therefore remove rivals' positive skill from
one comparison and increase its owner's normalized share through sibling
UIDs scoring elsewhere. This economic risk remains unresolved. The
bounded attack study in Section~\ref{development-and-launch-criteria}
must measure its frequency, collateral loss and owner-level payoff
before live rewards.

The versioned adjudication record binds the task and comparison IDs,
source/candidate/reference/brief commitments, evaluator and calibration
hashes, exact \(q_x,q_z,V\), rubric fields, evidence root, validator
snapshot and \(W\), ballot commitments/openings, field certificates,
deadlines, \(G\), and any unresolved or void reason. Records use the
canonical encoding and authenticated encrypted evidence paths of
Section~\ref{service-evidence}. Evidence must reach every authorized
certifier within the funded publication limits. Customer jobs never
enter this process automatically. A execution records additionally bind
the complete published model manifest and canonical runtime configuration.
Certificates attest scoped judgments and faithful computation. They do
not prove semantic correctness or require reproduction of private B
generation. A artifact loading, numerical replay, and full-matrix
capacity require their own implementation checks beyond the one-ticket
A-requester/B-provider rewrite mailbox pilot.
